\hypertarget{what-published-evaluations-establish-and-how-to-reuse-them}{%
\section{What published evaluations establish and how to reuse
them}\label{what-published-evaluations-establish-and-how-to-reuse-them}}

\label{exec14:evaluation-development}

\hypertarget{the-question-a-bank-must-answer}{%
\subsection{The question a bank must
answer}\label{the-question-a-bank-must-answer}}

A useful causal assistant must improve the decision made with the
available evidence at an acceptable total cost. Fluent reports,
agreement among agents and successful tool calls are intermediate
outcomes. The proposal therefore evaluates the target, assumptions,
graph, identification decision, numerical result and action separately.
A correct refusal can be valuable; an unsupported but precise number can
be harmful. This chapter develops a prospective comparison using public
evidence and stipulated mechanisms. It reports no new live-model
experiment, analyst-time observation or bank return on investment.

Consider the decision to offer a restructuring review at an account's
first qualifying arrears date. The target is the mean effect of the
offer on discounted 180-day net recovery among accounts eligible at that
date. The action is an offer, not eventual acceptance; information must
have been available before the offer, and the offer cost is counted
once. The experiment tests whether a workflow can translate that brief,
preserve the clock, decide identification and return an appropriate
estimate or refusal. It does not test whether the bank should deploy the
offer merely because the assistant completed an analysis.

\hypertarget{what-primary-evaluations-actually-establish}{%
\subsection{What primary evaluations actually
establish}\label{what-primary-evaluations-actually-establish}}

\emph{CLadder} by Jin et al.~\citep{v12_la_2312_04350}
(arXiv:2312.04350v3, Table 2) offers over 10,000 formal causal questions
and deliberately varies commonsense alignment. Its reported historical
GPT-4 accuracy rises from 62.03\% to 70.40\% with CausalCoT, an 8.37
percentage-point improvement. This is favorable evidence for structuring
a language-to-formal-query task. It is not evidence of discovering an
unknown bank graph or calibrating a treatment-effect interval. The
benchmark's formal answers are valuable only when the generator's
assumptions actually identify the named target.

\emph{CausalEval} by Yu et al.~\citep{exec13_census_causaleval}
(arXiv:2410.16676v4, Table 3 and Appendix C) supplies an important
counterweight. On its cost-limited GPT-4o subsets, function calling
scores 76 on CLADDER against 69 for CoT, while e-CARE falls from 90 to
73. These are source-reported percentages, not matched bank performance.
Its nine benchmarks have heterogeneous tasks; some models use 100
examples, and its human row combines previously published human scores
with three trained graduate students on selected datasets. Neither a
universal enhancement ranking nor a matched analyst productivity
comparison follows. The proposal must report results by task and
preserve a numerical, ordinary analyst and single-assistant baseline.

\emph{CausalDS} by Leban and Sun~\citep{v12_la_2607_08093}
(arXiv:2607.08093v1, Section 5, Figure 3) brings the evaluation closer
to an actual workflow: narrated scenes, files, a distinct observation
layer, private truth and explicit nonidentification. Its 100-scene exam
is sampled from 953 scenes, with 28 Rung-1, 51 Rung-2 and 21 Rung-3
tasks and 48/32/20 clean/proxy/proxy-hard conditions. The paper reports
substantially different content, abstention, interval and token
outcomes. However, the 20.0\%--71.4\% interval coverage figures use only
seven eligible ATE interval tasks per model, five for Qwen. They cannot
support a precise population calibration conclusion. Reasoning
effort/settings differ across model groups, and combined
prompt-plus-completion tokens are not dollar costs. The reported
failures from excessive tool output also show why bounded output belongs
in the experimental execution contract.

The study by Feng et al.~\citep{v12_la_2407_19638}, \emph{On the
Reliability of Large Language Models for Causal Discovery}
(arXiv:2407.19638v2, Sections 2--4) examines accessible OLMo/BLOOM
training corpora, manipulated relation frequencies and contextual
variants. Exposure and context affect responses. These findings motivate
separate tests of unfamiliar names, new mechanism families and
misleading evidence. Renaming a familiar graph does not establish
mechanism novelty. Equally, a changed context can legitimately change
the causal claim, so a context-sensitive answer is not automatically an
error. An independent target-specific adjudication is needed.

\emph{CausalChat} by Zhang et al.~\citep{v12_la_2410_14146}
(arXiv:2410.14146v1, Section 6) is evidence for useful interaction
design and a warning about productivity claims. Six university students
move through BIC feedback, LLM text and the interface in a fixed order,
choosing different remaining edges at each stage. The reported edge
errors are 5/18, 7/18 and 1/18. Learning, order and edge difficulty are
not separated by that design. Its expert participants also raise missing
citations. The proposal adopts traceable graph-edit suggestions and
measures review effort; it does not infer a causal labor-saving effect
from these counts.

The study by Le et al.~\citep{v12_la_2407_15073}, \emph{Multi-Agent
Causal Discovery Using Large Language Models} (arXiv:2407.15073v4,
ablations and limitations) reports a 16--25-fold token increase over
standard baselines, with coding/debugging taking up to half of the full
system's tokens. A single-agent variant is competitive on the simpler
graphs at approximately half the cost; full debate improves selected
harder examples. The study covers five public graphs with 5--20 nodes
and explicitly excludes hidden confounders and cycles. It provides a
reason to test debate against a single assistant, not to make debate
mandatory in every task or to infer its value for latent temporal
discovery.

Murphy's study~\citep{v12_la_2608_09696}, \emph{Model
Discovery Agent} (arXiv:2608.09696v5, Figure 2 and discussion) combines
language proposals with numerical fitting, simulation and experimental
design. Its chemistry comparison retains 115 matched evaluable
world--seed cells from 171 attempted, spanning 54 worlds. A conditional
performance curve for those cells is not all-attempt availability. The
comparisons are not compute-matched, its experiments use open-loop
policies, and the paper explicitly says finite held-out predictive tests
do not establish causal identification. The banking design retains
attempt counts and numerical computation costs and separates prediction,
experiment choice and effect identification.

Taken together, these studies support a testable architecture: typed
causal tasks, independent formal/numerical checks, source-grounded human
review and selective assistance. They do not establish one best model,
general causal competence or positive bank ROI. The reported results
remain literature evidence from distinct tasks, populations and model
settings.

\hypertarget{reusing-benchmarks-without-inheriting-their-labels-blindly}{%
\subsection{Reusing benchmarks without inheriting their labels
blindly}\label{reusing-benchmarks-without-inheriting-their-labels-blindly}}

The inspected CLadder source is pinned to commit \texttt{3d2d1169b4b9}
(2024-05-29), with an MIT notice. Its query generator contains a
Wald-ratio IV template and uses ATE bounds in constructing answers. An
unrestricted IV graph does not identify a population ATE through the
Wald ratio. Additional assumptions may identify a different target or a
sign; the full generator and item assumptions must be checked before
admitting a label. Quarantine IV-as-population-ATE cases and resolve
repository/paper dataset-version differences. Retain the generator as a
candidate source of formal fixtures after that review.

The inspected CausalDS source is pinned to \texttt{c74671d0a092}
(2026-07-17), with Apache-2.0 code and a CC0 data notice. Its separation
of conceptual SCM, observed proxy, private truth and scorer is useful.
The inspected grader can default missing identification metadata toward
identifiable and distinguishes intentional abstention from malformed
output. It computes both uncapped and capped normalized interval-score
summaries; publish both. Require explicit
identified/nonidentified/unknown states, finite endpoint validation,
preserved observation semantics and raw tail losses before reuse. These
are static source findings, not installed-package results or a completed
dependency/license audit.

The first prototype should import reviewed fixture definitions and
expected outputs through a narrow adapter, then compare the adapter
against an independently written oracle on tiny cases. The trusted
scorer must remain outside the agent's writable task environment. Freeze
hashes of code, data, prompts, retrieval corpus and scoring rules.
Impose tool-output and runtime limits while retaining the failure status
and incurred cost. Exposed development examples are useful for debugging
but cannot also be the hidden quality test. Neither upstream package has
been executed for this synthesis.~\citep{exec14_evaluation_code}

\hypertarget{a-reproducible-evaluation-example}{%
\subsection{A reproducible evaluation
example}\label{a-reproducible-evaluation-example}}

For an exact illustration, stipulate a population in which an offer is
randomized with probability one half, independent of binary pre-offer
risk \(X\) and potential outcomes. Let \(P(X=1)=1/2\) and let the
potential recovery means be \(E[Y(0)\mid X]=100+20X\) and
\(E[Y(1)\mid X]=105+30X\). The conditional effects are 5 and 15 and the
population ATE is 10 monetary units. Consistency, no interference and
the stipulated assignment/observation law identify it. Standardization
over the two known strata is an ordinary numerical baseline. The exact
means supply oracle truth; finite-sample intervals would require a
separately specified sampling law and method. In an observational
variant with an unmeasured common cause of offer and recovery and no
additional identifying restrictions, the appropriate answer is
nonidentification rather than the randomized formula.

A second, deliberately hypothetical accounting fixture has 100 requested
analyses: 80 identified requests and 20 nonidentified requests. Among
the identified requests, 70 return admissible intervals, 63 cover the
true target and 10 fail to return an interval. Among the nonidentified
requests, 16 correctly refuse and four return unsupported numbers. Then

\[
 \begin{gathered}
 A_{\mathrm{id}}=70/80=87.5\%,\\
 C_{\mathrm{returned,id}}=63/70=90\%,\\
 C_{\mathrm{return\ and\ cover,id}}=63/80=78.75\%.
\end{gathered}
\]

The correct-refusal rate is \(16/20=80\%\). Across all requests, the
fraction returning an admissible covering interval is \(63/100=63\%\),
but its complement is not an error rate: principled refusals occupy part
of it. The four unsupported numbers do not improve scientific
availability. An overall operational completion count must therefore sit
beside identification-specific outcomes and failure reasons, not replace
them. None of these stipulated counts estimates actual model
performance.~\citep{exec14_evaluation_oracle}

For a finite interval \([l,u]\) at nominal level \(1-\alpha\), the
interval score is \[
 (u-l)+\frac{2}{\alpha}(l-\theta)\mathbf1\{\theta<l\}
       +\frac{2}{\alpha}(\theta-u)\mathbf1\{\theta>u\}.
\] Report raw scores, coverage, width and availability with their
eligible denominators. A normalized score additionally needs a fixed
positive scale. Invalid endpoints are failures, not zero loss. A
disconnected identified set is evaluated as a set, retaining topology
and infinite length; do not replace it with a convenient midpoint or
force this finite-interval score onto it. Graph scores likewise use the
declared identifiable equivalence class; unsupported orientation is not
rewarded simply for resembling a fully directed reference graph.
